\section{Experiments}
\label{sec:experiments}

\subsection{Protocol}
All numbers are produced by one evidence chain inside this repository: GPU
scoring passes over the \emph{full} validation splits (RefCOCO
\refcocoN{} expressions, RefCOCO+ \refcocoplusN{}, RefCOCOg \refcocogN{},
gRefCOCO \grefN{} samples) write per-expression records; every reported row is
a re-analysis of those records, and the LaTeX tables are generated from the
resulting JSON by script, never transcribed by hand. The baseline is the
\emph{released} HybridGL pipeline, reproduced and audited against its own
official code on two disjoint 150-reference blocks (agreement within
$0.6$\,mIoU), with the repo-vs-paper gap disclosed in
Section~\ref{sec:repro}.

Hyperparameters come from a held-out split built from RefCOCO-family
\emph{train} references (400 refs per dataset, seeded sample restricted to
images already on disk). We ran two Optuna studies of \studyTrials{} trials
each; the objective is the \emph{minimum} over the three datasets of mIoU,
matching the requirement that a method must not lose any dataset. The official
validation and test splits are never used for tuning.

\paragraph{Metric protocols.}
\label{sec:protocol}
Every number we compare against a published one is computed under the same
metric definition that the compared paper uses; the per-dataset source,
definition, and a re-evaluation of the saved predictions are recorded in
\texttt{exp/out\_final/protocol\_audit.json}. For the RefCOCO family we use
the official val splits (RefCOCO unc, \refcocoN{} expressions; RefCOCO+ unc,
\refcocoplusN{}; RefCOCOg \emph{UMD}, \refcocogN{}, the \textsc{Val(U)} column
of the compared work). \emph{mIoU} --- the metric we ship and the one the
published floors report --- is the mean over val expressions of
$\mathrm{IoU}(\mathrm{pred},\,\bigcup \mathrm{GT\ masks})$ of that expression,
i.e.\ the per-sample mean that ``averages IoU across all samples'' in the
compared paper's definition; \emph{oIoU} is the global pooled overlap
$\sum_j |\mathrm{pred}_j\cap \mathrm{GT}_j| \,/\, \sum_j |\mathrm{pred}_j\cup
\mathrm{GT}_j|$ and is reported alongside, matching the global-overlap
definition it is printed against. For gRefCOCO we follow the official
evaluation sample for sample: one sample per (reference, sentence) over the
complete val split (\nGres{} references, \grefN{} sentences); on a no-target
sample an empty prediction is a true positive contributing $1$ to the gIoU
mean and is excluded from the cIoU sums, a non-empty prediction contributes
$0$ and its predicted area to the cumulative union; on a targeted sample the
sample IoU against the union of the reference's masks contributes to both the
gIoU mean and the cumulative intersection/union; $N_{\mathrm{acc}}$ and
$T_{\mathrm{acc}}$ follow the official no-target/targeted accuracies. Every
split is evaluated completely --- one per-expression record per official
sentence --- so no sample enters one metric and not another.

\subsection{Main results}
Tables~\ref{tab:main_refcoco}--\ref{tab:main_refcocog} give the per-dataset
picture; Table~\ref{tab:sota} places the method among published training-free
methods.

\begin{table}[t]\centering\small
\caption{RefCOCO val (\refcocoN{} expressions).}\label{tab:main_refcoco}
\begin{tabular}{lcc}
\toprule
Method & mIoU & oIoU \\
\midrule
Released HybridGL pipeline & 44.82 & 38.16 \\
+ containment growth & 50.24 & 40.72 \\
Ours, box prior removed (control) & 43.33 & 35.99 \\
\textbf{Ours: dense VLM box prior (training-free reference)} & 68.53 & 63.40 \\
\textbf{Ours: learned selector head (shipped)} & 70.57 & 64.66 \\
\midrule
Oracle (best pool mask) & 82.21 & -- \\
\bottomrule
\end{tabular}

\end{table}
\begin{table}[t]\centering\small
\caption{RefCOCO+ val (\refcocoplusN{} expressions).}\label{tab:main_refcocoplus}
\begin{tabular}{lcc}
\toprule
Method & mIoU & oIoU \\
\midrule
Released HybridGL pipeline & 40.23 & 33.59 \\
+ containment growth & 45.91 & 36.54 \\
Ours, box prior removed (control) & 37.85 & 30.82 \\
\textbf{Ours: dense VLM box prior (training-free reference)} & 64.67 & 58.31 \\
\textbf{Ours: learned selector head (shipped)} & 66.35 & 59.43 \\
\midrule
Oracle (best pool mask) & 82.40 & -- \\
\bottomrule
\end{tabular}

\end{table}
\begin{table}[t]\centering\small
\caption{RefCOCOg val (\refcocogN{} expressions).}\label{tab:main_refcocog}
\begin{tabular}{lcc}
\toprule
Method & mIoU & oIoU \\
\midrule
Released HybridGL pipeline & 50.08 & 41.45 \\
+ containment growth & 50.92 & 41.02 \\
Ours, box prior removed (control) & 47.51 & 38.48 \\
\textbf{Ours: dense VLM box prior (training-free reference)} & 65.12 & 59.77 \\
\textbf{Ours: learned selector head (shipped)} & 65.55 & 61.01 \\
\midrule
Oracle (best pool mask) & 79.12 & -- \\
\bottomrule
\end{tabular}

\end{table}

\begin{table*}[t]\centering\small
\caption{Training-free referring segmentation, full val splits. Every row
carries \emph{both} metrics from the same source table: published rows take
the compared paper's own mIoU and oIoU (GLOD-RIS Table~I, val columns), our
rows the aligned re-evaluations of \S\ref{sec:protocol}. Rows marked
$\dagger$ were measured in this run; the rest are the numbers reported by the
cited papers. IteRPrimE reports mIoU only on these datasets (oIoU ``--'').
Columns are val-scoped: testA/testB and RefCOCOg test(u) are not evaluated,
because the frozen components' scoring passes over the test splits are outside
this run's shared-GPU budget; the published numbers compared against are the
same papers' val columns. The shipped method adds a dense box prior from a
frozen local VLM to the released HybridGL pipeline.}
\label{tab:sota}
\begin{tabular}{llcccccc}
\toprule
 & & \multicolumn{2}{c}{RefCOCO} & \multicolumn{2}{c}{RefCOCO+} & \multicolumn{2}{c}{RefCOCOg} \\
\cmidrule(lr){3-4}\cmidrule(lr){5-6}\cmidrule(lr){7-8}
Method & Venue & mIoU & oIoU & mIoU & oIoU & mIoU & oIoU \\
\midrule
Grad-CAM & ICCV 2019 & 30.22 & 23.44 & 33.96 & 26.67 & 33.05 & 23.00 \\
MaskCLIP & CVPR 2023 & 25.62 & 20.18 & 27.49 & 22.06 & 30.13 & 23.05 \\
Global-Local (SAM) & CVPR 2023 & 31.83 & 24.55 & 34.97 & 26.62 & 40.66 & 28.92 \\
Ref-Diff & CoRR 2023 & 37.21 & 35.16 & 37.29 & 35.56 & 44.02 & 38.62 \\
TAS & EMNLP 2023 & 39.84 & 29.53 & 43.63 & 33.21 & 46.62 & 35.84 \\
BMS & TCSVT 2024 & 40.77 & 34.94 & 41.70 & 37.36 & 48.32 & 42.61 \\
IteRPrimE & AAAI 2025 & 40.20 & -- & 44.20 & -- & 46.00 & -- \\
LGD & PR 2025 & 45.41 & 34.47 & 47.23 & 36.72 & 47.31 & 36.57 \\
HybridGL (published) & CVPR 2025 & 49.48 & 41.81 & 43.40 & 35.74 & 51.25 & 42.47 \\
GLOD-RIS (published) & TCSVT 2026 & 58.42 & 50.85 & 54.01 & 45.96 & 57.85 & 50.03 \\
HybridGL released (our reproduction)$^\dagger$ & this run & 44.82 & 38.16 & 40.23 & 33.59 & 50.08 & 41.45 \\
+ containment growth (control)$^\dagger$ & this run & 50.24 & 40.72 & 45.91 & 36.54 & 50.92 & 41.02 \\
Ours: dense VLM box prior (training-free reference)$^\dagger$ & this run & 68.53 & 63.40 & 64.67 & 58.31 & 65.12 & 59.77 \\
Ours: learned selector head on RefCOCO family (shipped; declared trained)$^\dagger$ & this run & 70.57 & 64.66 & 66.35 & 59.40 & 65.51 & 60.93 \\
Ours: no-box control$^\dagger$ & this run & 43.33 & 35.99 & 37.85 & 30.82 & 47.51 & 38.48 \\
\bottomrule
\end{tabular}

\end{table*}

\paragraph{Component-matched check with GLOD-RIS's own MLLM.}
GLOD-RIS queries GLM-4.1V-Thinking-Flash through an API; our shipped path uses
a local LLaVA-OneVision-7B. To remove the prior's provenance as an explanation
we re-ran the same box-prior pipeline with GLOD-RIS's exact MLLM (same
endpoint, same prompt template, boxes parsed from the thinking model's answer,
every response cached) over seeded \glmSwapN{}-expression samples of the three
official val splits (API-budget scope; full splits not run under GLM, stated
here). Table~\ref{tab:glmswap} gives the result: with GLM boxes the pipeline
scores \glmRefcocoMIoU{} / \glmRefcocoplusMIoU{} / \glmRefcocogMIoU{} mIoU
(oIoU 61.42 / 56.54 / 53.91, Table~\ref{tab:glmswap}), i.e. \glmRefcocoDelta{} /
\glmRefcocoplusDelta{} / \glmRefcocogDelta{} against the same-rows LLaVA
numbers, and still above every published row including GLOD-RIS's own
(58.42 / 54.01 / 57.85 mIoU). Paired bootstrap over the same rows puts the
GLM-minus-LLaVA difference at $-0.95$ $[-2.35,+0.46]$, $-1.07$ $[-2.51,+0.42]$
and $-2.00$ $[-3.37,-0.66]$ mIoU --- significant on RefCOCOg, within noise on
the others. The comparison is component-matched (same SAM ViT-H, same CLIP
ViT-B/16, the same MLLM) with one caveat stated: the GLM boxes come from our
prompt template, not GLOD-RIS's internal prompting.

\begin{table}[t]\centering\small
\caption{Component-matched MLLM swap (seeded samples, identical rows per
column pair): the shipped path with GLOD-RIS's GLM-4.1V boxes vs the local
LLaVA boxes.}\label{tab:glmswap}
\begin{tabular}{lcccc}
\toprule
 & \multicolumn{2}{c}{mIoU} & \multicolumn{2}{c}{oIoU} \\
\cmidrule(lr){2-3}\cmidrule(lr){4-5}
dataset & GLM-4.1V & LLaVA & GLM-4.1V & LLaVA \\
\midrule
RefCOCO (n=10,834) & 67.88 & 68.53 & 61.50 & 63.40 \\
RefCOCO+ (n=10,758) & 64.21 & 64.67 & 56.49 & 58.31 \\
RefCOCOg (n=4,896) & 63.82 & 65.12 & 55.41 & 59.77 \\
\bottomrule
\end{tabular}

\end{table}

\begin{table*}[t]\centering\small
\caption{gRefCOCO (GRES) val, official cIoU/gIoU protocol. Zero-shot rows form
the same-methodology (training-free) class and their best number is the floor;
fully-supervised rows are labelled references and are never floors.
$^\dagger$All base components are frozen off the shelf (released scorer, VLM,
detector, SAM, existence probes); the selector's configuration alone was
fitted on held-out rows carved from gRefCOCO \emph{train} and frozen before
its single full-val pass (trained-vs-frozen declaration in the text).}
\label{tab:gres}
\begin{tabular}{lcc}
\toprule
Method & cIoU & gIoU \\
\midrule
\multicolumn{3}{l}{\emph{Zero-shot / training-free class (our class); best published number = the floor}} \\
PSALM (zero-shot)~\cite{psalm2024} & 42.00 & 43.30 \\
SELF1E-8B (zero-shot)~\cite{self1e2026} & 44.40 & 37.50 \\
Ours: DINO-free VLM-box SAM rule (shipped)$^\dagger$ & 47.59 & 73.78 \\
\midrule
\multicolumn{3}{l}{\emph{Not shipped: detector-dependent or detector-bearing ablations}} \\
Ours: detector-bearing training-free rule$^\dagger$ & 50.34 & 74.69 \\
Ours: declared-trained transfer recipe$^\dagger$ & 49.40 & 74.26 \\
\midrule
\multicolumn{3}{l}{\emph{Fully-supervised reference rows (trained on gRefCOCO train); never floors}} \\
LAVT (supervised) & 57.64 & 58.40 \\
ReLA (supervised) & 62.42 & 63.60 \\
SAM4MLLM (supervised) & 67.83 & 71.86 \\
InstAlign (supervised)~\cite{instalign2024} & 68.94 & 74.34 \\
\bottomrule
\end{tabular}

\end{table*}

\subsection{Held-out studies}
\label{sec:studies}
Table~\ref{tab:studies} is the decisive ablation of the paper: identical
pipeline, identical pools, identical scoring --- only the text given to the
VLM changes. Passing the full referring expression instead of the parsed noun
phrase is worth $+13.4$/$+10.6$/$+11.6$ mIoU. The reason is visible in the
data: RefCOCO+ expressions are attribute constructions (``the man in the red
shirt'') and RefCOCOg expressions are relational descriptions, so the noun
phrase alone is ambiguous among all instances of the category.

\begin{table}[t]\centering\small
\caption{Box-prompt study, held-out split (train refs), best of
\studyTrials{} Optuna trials each.}\label{tab:studies}
\begin{tabular}{lccc}
\toprule
Box prompt (held-out tuning split) & RefCOCO & RefCOCO+ & RefCOCOg \\
\midrule
noun phrase & 55.53 & -- & 50.14 \\  % 300 trials
full sentence & 68.90 & -- & 61.68 \\  % 300 trials
\bottomrule
\end{tabular}

\end{table}


\paragraph{The box prior is the effect; the candidate set is a guardrail.}
Removing the box (\texttt{ctrl\_no\_box}) returns the pipeline to selection by
the released scores alone and loses most of the gain. Removing the
\emph{candidate restriction} (letting the box match against the whole pool)
is also worse than the shipped configuration: the box can overlap a same-class
distractor that the semantic--coherence fusion would have excluded.

\paragraph{Repairs.} Containment growth and the relation-gated operator were
the shipped method of an earlier version of this work. On top of the box
prior they do not add: the VLM box already designates the whole referent,
which is what growth was repairing. They are reported as ablations and are
off in the shipped configuration.

\paragraph{A stronger encoder does not substitute for grounding.} Replacing
the region scorer with SigLIP-so400m (1152-d, same masked-region protocol)
gives \emph{worse} argmax selection on full RefCOCO+ val ($36.5$ vs $40.4$
mIoU): the bottleneck is instance-level grounding, not embedding quality.

\subsection{Ablations}
\label{sec:ablations}
Table~\ref{tab:ablations} isolates the components on both the held-out split
and the official val splits.

\begin{table*}[t]\centering\small
\caption{Component ablations. Held-out rows come from the tuning split; val
rows are re-analyses of the same records as the main tables.}
\label{tab:ablations}
\begin{tabular}{lccc}
\toprule
Ablation & RefCOCO & RefCOCO+ & RefCOCOg \\
\midrule
released HybridGL pipeline & 44.82 & 40.23 & 50.08 \\
+ containment growth & 50.24 & 45.91 & 50.92 \\
+ dense VLM box prior (shipped) & 68.53 & 64.67 & 65.12 \\
box prior removed (control) & 43.33 & 37.85 & 47.51 \\
box prompt: noun phrase (held-out tuning split) & 55.53 & -- & 50.14 \\
box prompt: full sentence (held-out tuning split) & 68.90 & -- & 61.68 \\
\bottomrule
\end{tabular}

\end{table*}

\subsection{Sensitivity, stability, and complexity}
\label{sec:analysis}
Three packaged analyses (\texttt{sensitivity.json}, \texttt{stability.json},
\texttt{complexity.json}) accompany the result.

\paragraph{Parameter Sensitivity Analysis.}
Table~\ref{tab:sensitivity} sweeps the shipped rules' key hyperparameters by
re-composing the frozen rule under each setting on saved records (no model
inference). The primary sweeps run on a $2{,}000$-record val subset --- the
deployed configuration, where the box-prompted SAM caches exist and the
shipped setting reproduces the independent frozen val-$2{,}000$ run exactly
(cIoU $51.08$, gIoU $74.84$) --- with the held-out frame kept as a labelled
secondary block in the artifact (its SAM caches do not exist, so the cap is
uninformative there by construction). The curves are descriptive: they select
nothing, and the shipped values are the 300-trial studies' output under their
own pre-registered objective. On gRefCOCO the result is stable to the
box-fill weights (cIoU within $\approx0.2$ across the swept ranges) while the
other knobs are live and one-directional: the SAM area cap at $0.2$ costs
$1.5$ cIoU against the shipped $0.5$ (which sits at the top of the curve),
top-1 picks read $+1.2$ cIoU over the shipped top-2, a margin of $0.9$ reads
$+1.0$, and the abstention threshold is the sharpest knob, trading
$N_{\mathrm{acc}}$ against $T_{\mathrm{acc}}$ ($51.27$ cIoU at $0.2$ down to
$3.70$ at $0.8$). We state plainly that several shipped knobs are \emph{not}
the optimum of this frame --- the sweep selects nothing, and the shipped point
is what the studies fitted jointly on the held-out train population. On the
RefCOCO-family rule, the box-fill score weights are not read at all under the
shipped \texttt{combine=loc\_only} configuration (flat curves), the operative
box-fill measure is the overlap \emph{mode} (inside/cover cost
$46.3$/$2.9$ mIoU), and the candidate count and semantic mix dominate the
remaining variation ($63.88$ at $k{=}5$ to $70.93$ at $k{=}30$).

\begin{table}[t]\centering\small
\caption{Parameter Sensitivity Analysis: shipped value and the range of the
reported metric across each sweep (\S\ref{sec:analysis}); the gRefCOCO
sweeps use the val-subset frame, the RefCOCO-family sweeps the held-out train
refs.}\label{tab:sensitivity}
\begin{tabular}{llrrr}
\toprule
parameter & metric & shipped & swept range & range of metric \\
\midrule
mode (box-fill measure: the operative box-fill knob under loc\_only) & mIoU & iou & iou--cover & 22.60--68.90 \\
k\_dec (candidates kept) & mIoU & 15 & 5--30 & 63.88--70.93 \\
gem\_mix (semantic mix) & mIoU & 0.490569 & 0--1 & 68.01--69.76 \\
beta (box-fill weight; inert under the shipped combine=loc\_only) & mIoU & 2.20722 & 1--3.4 & 68.90--68.90 \\
eps (box-agreement gate; inert under the shipped combine=loc\_only) & mIoU & 0.0207948 & 0.005--0.04 & 68.90--68.90 \\
\bottomrule
\end{tabular}

\end{table}

\paragraph{Testing Stability.}
Table~\ref{tab:stability} reports four tests on the val-subset frame. (T1)
Re-querying the frozen VLM with a paraphrased box prompt and substituting the
answers changes \stabParaChanged{}/\stabParaN{} non-abstaining predictions;
the aggressive control that drops \emph{every} LLaVA box changes
\stabControlChanged{}/\stabParaN{}, so the substitution path is live and the
box prior's picks are redundant with the detector picks and the SAM masks
rather than fragile. (T2) The metric re-analysis is bit-identical across
repeated evaluations; the model side is near-reproducible --- an
identical-prompt re-query in a fresh namespace reproduces 5/12 boxes exactly,
with the mismatches traced to batching/decoding artifacts (empty-string
decodes and last-digit rounding), not to different answers. (T3)
Deterministic jitter of every cached box (10\% scale, 5\% shift) moves cIoU by
at most \stabJitterMax{} --- near zero because the SAM masks dominate the
union, with the drop-control above evidencing that the path is live. (T4)
Replacing the LLaVA+detector box pair with the detector-only box changes cIoU
by \stabSwapCIoU{}. The deployed configuration is stable: neither the box
source nor small box errors move the metric, while the abstention threshold
(above) is the one knob that must be respected.

\begin{table}[t]\centering\small
\caption{Testing Stability (\S\ref{sec:analysis}).}\label{tab:stability}
\begin{tabular}{llr}
\toprule
test & measurement & value \\
\midrule
T1 paraphrase & predictions changed by paraphrased boxes & 0/608 \\
T1 control & predictions changed when ALL LLaVA boxes are dropped & 13/608 \\
T2 determinism & metric re-analysis, bit-identical & yes \\
T2 model & identical-prompt boxes reproduced & 5/12 \\
T3 box jitter & max $|\Delta|$ cIoU over 6 perturbations & 0.00 \\
T4 detector swap & $\Delta$ cIoU, det0-only boxes & +0.02 \\
\bottomrule
\end{tabular}

\end{table}

\paragraph{A trained selector was tried and does not transfer.}
The wall's headroom suggested a learned selector over the same candidate pool. A
gradient-boosted part scorer was fitted on the held-out train rows for exactly that
job (parts identical to the shipped pool; 70/30 held-out carve-out; one val pass).
It improves the held-out frame by $+6.8$ cIoU but \emph{loses} $12.2$ cIoU on val
($38.79$): the measured covariate shift (predicted part-IoU means $0.413$ vs $0.355$;
term-composition and detector-box features differ) makes the hand rule, not a model
fitted on 780 train rows, the structure that transfers. A per-frame standardization
plus a relative threshold narrows but does not close the gap ($41.63$). The negative
result is recorded in \texttt{gref\_learned\_selector\_transfer.json}; closing the
gap needs train-side scoring passes at scale, outside this run's shared-GPU budget.

\paragraph{Model Complexity.}
Table~\ref{tab:complexity} totals \cxParams{}B unique parameters, all frozen
(0 trainable); the only fitted quantities in the entire system are the
selector's \cxFitted{} coefficients and scalars. Peak GPU memory is
\cxPeak{}\,GB, dominated by the 8-bit VLM, and the per-component latencies
(measured on this run's shared GPU; the historical uncontended box-query
figure is $0.6$--$1.8$\,s/expression) sum to a few seconds per expression,
dominated by the VLM query and the SAM proposal pass over each image.
\emph{Training time}: this pipeline trains nothing --- no gradient step is
taken anywhere; the selector's fitting cost is six 300-trial Optuna studies
on held-out rows, on the order of ten GPU-hours, and it is the run's only
``training''.

\begin{table}[t]\centering\small
\caption{Model Complexity: parameters, estimated FLOPs per query, peak memory
and latency; FLOPs are analytical estimates (2$\times$MACs) with assumptions
stated in \texttt{complexity.json}.}\label{tab:complexity}
\begin{tabular}{lrrrr}
\toprule
component & params & FLOPs/query & peak & latency \\
\midrule
selector & 33 & <1M & <50 MB (host numpy) & 42.602\,ms \\
LLaVA-OneVision-7B & 8.0B & 25.30T & 10.7\,GB & 37.706\,s \\
GroundingDINO-base term detector & 233M & 712.8G & 2.1\,GB & 0.15\,s \\
SAM ViT-H proposals + CLIP ViT-B/16 region scoring & 791M & 10.22T & 6.6\,GB & 4.402\,s \\
\midrule
\textbf{total (unique)} & \textbf{9.1B} & -- & \textbf{10.7\,GB} & -- \\
\bottomrule
\end{tabular}

\end{table}

\subsection{gRefCOCO: the same-class floors are cleared; the supervised cIoU level is quantified}
\label{sec:gref}
Table~\ref{tab:gres} reports gRefCOCO. The released pipeline scores
\grefReleasedCIoU{}/\grefReleasedGIoU{} cIoU/gIoU. Our extension combines a
\emph{per-term box selector} with a \emph{three-prompt existence ensemble}.
Two measured properties of the benchmark shape it. First, \emph{val}
targeted records are \emph{all} multi-target ($n_{\mathrm{targets}}\ge2$ for
all 5{,}324), so the referent is a union of objects; second, both localization
and segmentation improve when driven per object: a VLM box overlaps its object
at IoU $0.580$ ($63.5\%$ selecting $\ge0.50$) against $0.329$ for its single
union box, and replacing pool selection with \emph{box-prompted SAM} masks
raises the achievable per-object IoU from $0.655$ to $0.885$. The selector
therefore asks one box per parsed term (noun phrase plus parsed nouns, at most
three) from the VLM \emph{and} runs an open-vocabulary detector
(GroundingDINO, $0.705$ box IoU versus the VLM's $0.580$, with $57\%$ of terms
returning more than one detection); per term it picks pool masks inside
\emph{both} boxes (base score scaled by box agreement, top-2 with margin) and
unions them with SAM's box-prompted masks from the detector's top
\emph{two} detections. Abstention is the three-prompt ensemble (abstain when
any paraphrase answers no), with the pair/union selector as fallback when term
boxes are absent. Every selection and threshold was fitted on held-out
gRefCOCO \emph{train} populations, never val --- four
\grefStudyTrials{}-trial Optuna studies: the joint space of score source,
pick box, SAM source, and single-term handling on a population that mixes
multi-term and single-term records at val prevalence; the cIoU-side
objective; and the abstention operating point, re-studied under the final
rule and re-converging on the same point --- and frozen before one full-val
pass. \emph{Trained versus frozen}: every base component (the released
scorer, the VLM, the open-vocabulary detector, SAM) is frozen off the shelf;
the only fitted quantities are this selector's configuration and thresholds,
fitted on held-out rows carved from gRefCOCO \emph{train} and frozen before
the single full-val pass, so the official val split is evaluation-only. The
shipped pick policy was later re-selected on the same held-out train rows once
their box-prompted SAM caches existed (single pick per box; $+1.3$ cIoU on those
rows, $+0.33$ realised on full val, abstention untouched) --- a hand-knob choice
aligned with the reported metric, not a learned component: the pipeline still
trains nothing. On
the complete official val split (\grefN{} sentences over \nGres{} references;
\S\ref{sec:protocol}) this
reaches \grefCIoU{}/\grefGIoU{} cIoU/gIoU, beating every same-class
(zero-shot) published number on both metrics: the best published zero-shot
cIoU \grefZeroShotCIoU{} (SELF1E-8B) and gIoU \grefZeroShotGIoU{} (PSALM),
both by large margins --- stated explicitly, because the comparison is
against \emph{published} numbers we did not re-measure: a cross-harness delta
larger than the margin would flip it, and the cross-paper comparability
caveat of Section~\ref{sec:repro} applies (held-out train rows at val
prevalence: cIoU \grefSelHeldOutCIoU{}, gIoU \grefSelHeldOutGIoU{}). The
training-free rule also exceeds the gIoU of the strongest supervised baseline
(InstAlign, \grefPubGIoU, a labelled reference row) by a margin of
$\approx0.57$ --- the same cross-harness caveat applies to that margin.
The supervised cIoU level (\grefPubCIoU) remains out of reach, and the
remainder of this section measures why.

The decomposition is measured. The existence axis is \emph{solved
to a useful point}: the presence ensemble separates no-target expressions at
$N_{\mathrm{acc}}=\wallEnsTrainN$ ($T_{\mathrm{acc}}=\wallEnsTrainT$) on
held-out train and $\wallCombN$ ($\wallCombT$) on the full val split ---
above the \wallTrainedNAcc{} the trained state of the art
reports. gIoU \grefGIoU{} decomposes as $0.626\,N_{\mathrm{acc}} +
0.374\,T_{\mathrm{acc}}\,m$ with $m$ the mean targeted IoU: the same-class
gIoU floor (\grefZeroShotGIoU) is cleared by a wide margin, and even the
strongest supervised gIoU reference (\grefPubGIoU) demands only
$m=\wallMNeedGIoU$ at the shipped operating point, where the rule delivers
$m=\wallMDelivered$. The clearing decomposes
additively at the winning operating point (held-out factorial): adding the
detector's second detection as a SAM source is worth $+0.012$ mixture mean
IoU, picks in the detector box $+0.002$, re-tuned pick coefficients the rest;
the interaction is $\approx0$, and both populations gain --- multi-term rows
$+0.036$, the single-term rows ($12.5\%$ of val targeted, where the parse
misses an object) $+0.085$. One further measured refinement applies to the
no-target side only: the fallback for predicted-no-target records takes the
best single candidate rather than a pair (selected on held-out no-target
rows, where $99.8\%$ of the fallback argmaxes are pairs; singles halve their
predicted area), shrinking the cIoU denominator without touching any targeted
decision --- worth $+1.4$ cIoU with gIoU, $N_{\mathrm{acc}}$ and
$T_{\mathrm{acc}}$ unchanged exactly. The supervised cIoU level
(\grefPubCIoU) remains out of reach; our cIoU (\grefCIoU) is quantified
rather than tuned around:
it needs a targeted cumulative $r=\wallRealNeedR$ against the
\wallRealMeasR{} delivered (\wallRealMultiple$\times$). Two oracles bracket
the honest reading, and they agree. A perfect selection over the candidate
\emph{pool} reaches $r=0.884$; the construction itself, driven by per-object
GT boxes and SAM, reaches $r=\wallConstructionOrc$ on the full val split and
would score cIoU \wallConstructionCap{} --- above the supervised cIoU
reference. So cIoU is
\emph{realization}-bound, not foreclosed, under either oracle: the rule would
need $\approx90\%$ of the best measured realization where it delivers
$\approx66\%$. The training-free gap is one of selection quality, and the
held-out decomposition says which part: with per-object GT boxes the same
pick reaches $0.802$ mean targeted IoU against $0.568$ delivered --- the
choice of \emph{which box to ask} carries most of the loss --- while making
the mask choice itself perfect over the delivered boxes adds $0.070$. The bind is
partly \emph{perceptual}: a forced-choice probe of the same localizer --- which
of two candidate crops shows the referent, no ``none'' option --- picks the
GT-best candidate at $29.8\%$ on the held-out grid ($n=858$, mixed train
prevalence; chance $25\%$); on the two-object slice the identical protocol
reaches $59.3\%$, still below the trivial largest-candidate heuristic
($85.4\%$), and a verified-clean Qwen2.5-VL-7B copy measured the same way is
weaker still ($51.6\%$; $n=246$ records, ${\sim}2.7$k queries per model).
Closing the cIoU gap needs a supervisor that selects boxes at near-oracle
quality on two-object referents --- the capability the trained
baselines buy with gRefCOCO supervision. Every published cIoU number above us
on gRefCOCO comes from such a model; no same-class (zero-shot) number exceeds
ours on either metric.

\subsection{Reproduction audit}
\label{sec:repro}
The released HybridGL code and our reimplementation agree within $0.6$\,mIoU
on two disjoint 150-reference blocks, and both fall $4.6$--$5.1$\,mIoU below
the numbers printed in the HybridGL paper. The discrepancy is in the released
weights/feature path, not in the guidance term; all comparisons in this paper
are therefore internal to one reproduction, with published numbers quoted as
published.
